Large language models can generate machine-readable recommendations for decomposing monolithic applications, but a valid output does not necessarily identify useful service boundaries. This paper presents TraceMicro, an evidence-linked workflow for assigning Java classes to candidate services, and examines its recorded behavior on a PetClinic development snapshot. The system combines source and static-dependency evidence, structured model responses, LangGraph orchestration, shared-class support, deterministic validation and bounded request accounting. We compare package grouping, a deterministic dependency baseline and three completed full-scope model executions over 25 eligible classes. The deterministic methods produce group sizes of 11/5/4/3/2 and 21/1/1/1/1, respectively. All three model executions satisfy coverage and citation-reference checks yet assign every class exclusively to one service. Together they consume 27 requests and 144,478 reported tokens. A separate three-class interactive pilot produces two overlapping groups, while two full-scope planning configurations fail before assignments begin. These observations demonstrate an inspectable execution pipeline and a repeatable single-group outcome. Architectural quality remains unmeasured. We distinguish software verification, operational feasibility and architectural assessment, and identify the independent references and controlled experiments needed to evaluate proposed extensions.
